\section*{Chapter \ref{ch:rs:from-signal-to-forecast} --- From Signal to Forecast}

\begin{solution}{exo:rs:from-signal-to-forecast:1}
$0.04 \times 10\% \times (-1.5) = -0.6\%$ for the month; $0.04 \times 5\% \times 2 = 0.4\%$.
\end{solution}

\begin{solution}{exo:rs:from-signal-to-forecast:2}
$0.03\sqrt{12 \times 500} = 2.32$; $0.03\sqrt{52 \times 50} = 1.53$, if the weekly forecasts are independent, which a slow signal's are not.
\end{solution}

\begin{solution}{exo:rs:from-signal-to-forecast:3}
2 and 1 violate the order: pool to 1.5, 1.5. 3, 3 are in order; the next 2 violates: pool the last 3 with it (2.5), which still violates the first 3:
pool the three to $8/3 = 2.667$. 4 is in order. Fit: 1.5, 1.5, 2.667, 2.667, 2.667, 4.
\end{solution}

\begin{solution}{exo:rs:from-signal-to-forecast:4}
$100/(1 + 99 \times 0.05) = 16.8$. To halve it to 8.4: $1 + 99\rho = 100/8.4$, $\rho = 0.110$.
\end{solution}

\begin{solution}{exo:rs:from-signal-to-forecast:5}
Qian and Hua: $0.04/0.05 \times \sqrt{12} = 2.77$. The law: $0.04\sqrt{12 \times 704} = 3.68$. The law assumes the IC varies only by sampling
across 704 independent names, a standard deviation of about $1/\sqrt{704} = 0.038$; the measured 0.05 includes the factors' common moves, which the
law ignores.
\end{solution}

\begin{solution}{exo:rs:from-signal-to-forecast:6}
Over one day, $0.02 \times 2\% = 0.04\%$. Over a month, the sum of the single-day ICs, $0.02 + 0.01 = 0.03$, times the daily volatility: $0.06\%$; not
21 times the daily forecast.
\end{solution}

\begin{solution}{exo:rs:from-signal-to-forecast:7}
\texttt{rs\_forecast.concentrated}: the \omterm{def:rs:from-signal-to-forecast:breadth}{transfer coefficient} falls to 0.59 and the information ratio to 2.18, against 0.89 and 3.25 around the equal-weight
benchmark. Most names now have tiny benchmark weights, so the book can underweight them only by their tiny weight, while the few large names absorb the
negative bets: the risk-adjusted weights lose their resemblance to the forecasts.
\end{solution}

\begin{solution}{exo:rs:from-signal-to-forecast:8}
\omterm{def:rs:from-signal-to-forecast:law}{Breadth} counts independent forecasts: a signal whose ranks change little from day to day (chapter~13) gives far fewer than 252 new forecasts a year,
and 3\,000 names that share industries and factors are fewer than 3\,000 independent bets. The daily IC of a slow signal is also smaller than its
monthly IC, the portfolio's constraints lower the \omterm{def:rs:from-signal-to-forecast:breadth}{transfer coefficient}, and trading daily costs money. The formula's 26 is a ceiling nobody reaches.
\end{solution}

\begin{solution}{pb:rs:from-signal-to-forecast:1}
\begin{enumerate}
\item $0.05 \times 7.7\% \times 2 = 0.77\%$ for the month.
\item 0.055.
\item Slope 0.92 and $R^2$ 0.24\% for the scaled forecast; slope 0.0039 on the raw score.
\item An IC of 0.05 explains about $0.05^2$ of the variance of each name's return; the forecast is still valuable because a portfolio sums hundreds of
such small edges.
\item When the relation is visibly non-linear in the tails or saturates; with enough data in the tails to fit steps reliably.
\item 4.60.
\item 4.92; the sampling error of an information ratio over ten years of monthly returns is about 0.45, so yes.
\item Market residuals: mean 0.0505, standard deviation 0.0424. Gaussian residuals: standard deviation 0.0331.
\item 569 independent forecasts a month, against 704 names.
\item 0.86, with 47\% of holdings clipped at zero.
\item 3.54 predicted; 3.00 realised.
\item From 0.98 and 3.91 at 0.25\% a month to 0.86 and 3.00 at 4\% (table).
\item \textbf{Named result.} Of the law's 4.60, the IC's noise on the market's residuals (569 effective names of 704) leaves 4.13, the long-only \omterm{def:rs:from-signal-to-forecast:breadth}{transfer
coefficient} of 0.86 leaves 3.54, and the constraint's noise leaves the realised 3.00: two thirds of the promise.
\item Lower it again, in proportion to turnover: the long-only book at a high requested tracking error trades the most.
\item Lower it (0.59 in exercise 7): small benchmark weights leave little room for negative bets.
\item More names add correlated bets (little \omterm{def:rs:from-signal-to-forecast:breadth}{effective breadth}), smaller and less liquid positions, and more estimation error in the risk model and the
IC.
\item The law's promise, the IC's mean over its standard deviation, the \omterm{def:rs:from-signal-to-forecast:breadth}{transfer coefficient} and the realised information ratio, side by side.
\item Research improves the IC and its stability, and the \omterm{def:rs:from-signal-to-forecast:breadth}{effective breadth} by finding independent signals; portfolio construction improves the \omterm{def:rs:from-signal-to-forecast:breadth}{transfer
coefficient} and the constraints' noise.
\item It promises an information ratio that ignores the IC's noise, the constraints and costs, and invites adding names or trading more often.
\item The information ratio that forecasts of a given quality could deliver if they were independent, unconstrained and free to trade.
\end{enumerate}
\end{solution}

\begin{solution}{iq:rs:from-signal-to-forecast:1}
$\mathrm{IR} \approx \mathrm{IC}\sqrt{\mathrm{BR}}$ (Grinold, 1989): the information ratio of the optimal portfolio grows with the forecasts' correlation with
realised returns and with the square root of the number of independent forecasts a year. Assumptions: independent bets, a constant IC, an unconstrained
mean-variance portfolio, no costs.
\end{solution}

\begin{solution}{iq:rs:from-signal-to-forecast:2}
Standardise the score by date, estimate its IC on past data, and set $\alpha_i = \mathrm{IC}\cdot\omega_i z_i$ with the risk model's residual volatility
over the horizon; check calibration out of sample (Mincer--Zarnowitz slope near 1, \omterm{def:rs:from-signal-to-forecast:curve}{calibration curve} by bins).
\end{solution}

\begin{solution}{iq:rs:from-signal-to-forecast:3}
The IC's noise (its standard deviation against $1/\sqrt n$: correlated bets), the \omterm{def:rs:from-signal-to-forecast:breadth}{transfer coefficient} of the constraints, the constraints' noise, costs,
and whether the IC used in the law was estimated in sample. Measure each: the chapter's planted forecast lost 10\%, 14\% and 15\% at the first three.
\end{solution}

\begin{solution}{iq:rs:from-signal-to-forecast:4}
The correlation between the portfolio's risk-adjusted active weights and the risk-adjusted forecasts. Long-only constraints (especially against concentrated
benchmarks), turnover limits, sector and factor neutrality, position limits and liquidity caps lower it.
\end{solution}

\begin{solution}{iq:rs:from-signal-to-forecast:5}
Out of sample: regress realised on forecast returns (intercept near 0, slope near 1), plot mean realised against mean forecast by forecast decile, compare by
volatility bucket and by period, and refit the calibration only on data before each date.
\end{solution}

\begin{solution}{iq:rs:from-signal-to-forecast:6}
With $x = w\omega$ normalised and $u = r/\omega$ with mean $\mathrm{IC}\,z$ and unit independent noise, the active return has mean $\mathrm{IC}\,x^\top z =
\mathrm{IC}\sqrt n\,\mathrm{TC}$ and unit risk; annualising over independent periods gives $\mathrm{TC}\cdot\mathrm{IC}\sqrt{\mathrm{BR}}$. The independence of the
noise fails first: common factors make bets move together, which is why the chapter's IC had a larger standard deviation on the market's residuals than on
Gaussian ones.
\end{solution}
